% !TeX root = ../main.tex

Das folgende Kapitel behandelt die Implementierung.
Am Anfang wird die Entwicklungsumgebung thematisiert, dabei wird sowohl auf die Hardware-, als auch auf die
Softwarekomponente geschaut.
Danach wird auf den Aufbau der Software eingegangen, indem die Projektstruktur näher betrachtet wird.
Weiterhin wird die Implementierung der Fehlererkennung erklärt sowie die Implementierung eines Ablaufs.
Dazu wird die Konfiguration und ihre Parameter erklärt.
Als Letztes werden Implementierungsentscheidungen genauer erläutert.

\section{Entwicklungsumgebung}
\label{sec:devenv}
% !TeX root = ../../main.tex

Die Experimente wurden auf einem Testsystem ohne Hintergrundlasten durchgeführt.
Um externe Störfaktoren zu minimieren, wurden sämtliche Netzwerk- und Internetverbindungen für die Dauer der
Testreihen getrennt.
Das Testsystem war mit einem AMD Ryzen 7 7800X3D (8-Core Processor), 32 GB DDR5-RAM sowie dem Betriebssystem
Windows 11 Home ausgestattet.

Für die Implementierung kam die Programmiersprache Python (Version 3.14) zum Einsatz.
Als Datenbankmanagementsystem wurde PostgreSQL (Version 16.13) genutzt, wobei die Anbindung über psycopg2 erfolgte.
Diese wurde aufgrund ihrer Threadsicherheit ausgewählt.
Die anschließende Datenanalyse wurde mit pandas durchgeführt, um eine performante und strukturierte Auswertung zu
gewährleisten.

Die Entwicklung erfolgte in JetBrains IntelliJ IDEA Ultimate und für das Versionsmanagement wurde GitLab eingesetzt.


\section{Aufbau der Software}
\label{sec:softwareDesign}
% !TeX root = ../../main.tex

Der strukturelle Aufbau des Systems ist in Abbildung~\ref{fig:softwareDesign} dargestellt und wird im
nachfolgenden Abschnitt detailliert beschrieben.
% !TeX root = ../main.tex

\begin{figure}[ht]
	\centering
	\begin{tikzpicture}
		\node[
			fill=TerminalBackground,
			draw=TerminalFrame,
			rounded corners,
			inner sep=10pt
		] {
			\begin{tikzpicture}[
				scale=0.8,
				transform shape,
				node distance=1cm,
				every node/.style={font=\small},
				startstop/.style={
					rectangle,
					rounded corners,
					draw,
					minimum width=3cm,
					minimum height=1cm,
					align=center
				},
				process/.style={
					rectangle,
					draw,
					minimum width=3cm,
					minimum height=1cm,
					align=center
				},
				decision/.style={
					diamond,
					draw,
					align=center,
					aspect=2
				},
				storage/.style={
					rectangle,
					draw,
					dashed,
					align=center,
					aspect=2
				},
				arrow/.style={
					->,
					>=stealth
				},
				implements/.style={
					dashed,
					-{Triangle[open]}
				},
			]

				\node (start) [startstop] {main.py\\Programmstart};
				\node (runner) [process, below=of start] {experimentRunner.py};
				\node (analysis) [process, right=2cmof runner] {analysis.py};
				\node (experiment) [process, below=of runner] {experimentFunctions.py};
				\node (config) [process, left=2cm of runner] {config.py};
				\node (database) [process, below=of config] {database.py};
				\node (retry) [process, below=of experiment] {retry.py};
				\node (service) [process, below=of database] {service.py};
				\node (log) [storage, below=of analysis] {Logs};

				\draw [arrow, shorten <=5pt, shorten >=5pt] (start) -- (runner);
				\draw [arrow, shorten <=5pt, shorten >=5pt] (start) -| (analysis);
				\draw [arrow, shorten <=5pt, shorten >=5pt] (runner) -- (experiment);
				\draw [implements, shorten <=5pt, shorten >=5pt] (config) -- (experiment);
				\draw [arrow, shorten <=5pt, shorten >=5pt] (experiment.south) to[bend right=20] (retry.north);
				\draw [arrow, shorten <=5pt, shorten >=5pt] (retry.north) to[bend right=20] (experiment.south);
				\draw [arrow, shorten <=5pt, shorten >=5pt] (retry.west) to[bend left=20] (service.east);
				\draw [arrow, shorten <=5pt, shorten >=5pt] (service.east) to[bend left=20] (retry.west);
				\draw [arrow, shorten <=5pt, shorten >=5pt] (service.north) to[bend left=20] (database.south);
				\draw [arrow, shorten <=5pt, shorten >=5pt] (database.south) to[bend left=20] (service.north);
				\draw [implements, shorten <=5pt, shorten >=5pt] (database) -- (experiment);
				\draw [implements, shorten <=5pt, shorten >=5pt] (log) -- (analysis);
				\draw [arrow, shorten <=5pt, shorten >=5pt] (experiment) -- (log);

			\end{tikzpicture}};
	\end{tikzpicture}

	\caption{Aufbau der Systemstruktur}
	\label{fig:softwareDesign}
\end{figure}

Der Programmablauf beginnt mit der Ausführung der \texttt{main}-Funktion.
Von dort aus werden die Funktionen aus dem Modul \texttt{experimentRunner} und die
\texttt{start\_analysis}-Funktion nacheinander aufgerufen.
Die Funktionen des Moduls \texttt{experimentRunner} dienen dabei ausschließlich als namensgebende Funktionen, aus denen
im weiteren Verlauf die verschiedenen Fehlerklassen abgeleitet werden.
Dadurch wird klar zugeordnet, welche Fehlerklassen erzeugt und durch die Experimente getestet werden.
Die anschließende Analyse der Experimente erfolgt über die im Modul \texttt{analysis} enthaltenen Funktionen,
welche die erzeugten Ergebnisse weiterverarbeiten, auswerten und strukturiert darstellen.

Die in \texttt{experimentRunner} implementierten Funktionen rufen die \texttt{start()}-Funktion im Modul
\texttt{experimentFunctions} auf.
Hier sind neben den Experimenten auch die Funktionen zum \emph{Warm-up} der Datenbank und zum Schreiben der Logs
implementiert.
Zusätzlich beziehen sich die Experimente die Parameter aus dem \texttt{config} Modul.
In diesem Modul werden die Clients gestartet, die den Retry ausführen und für jeden Client wird eine neue
Datenbankverbindung aufgebaut.
Die Clients starten daraufhin den Retry mit einer bestimmten Strategie.

Im \texttt{retry}-Modul befinden sich die Implementierungen der einzelnen Retry-Strategien.
In jeder Strategie wird als Erstes die Transaktion ausgeführt.
Dafür ruft eine Funktion den Service auf, der zu der Fehlerklasse passt, den die Funktion aus
\texttt{experimentRunner} am Anfang definiert hat.
In dem Service steht die Implementierung der Transaktion, die den Fehler erzeugen soll.
Diese wird dann von der Datenbank ausgeführt und diese gibt dann entweder den SQLSTATE-Code \texttt{00000} oder den
Fehlerklassen spezifischen Fehler-Code.
Diese werden dann in \texttt{retry} ausgewertet und gespeichert.
Daraufhin setzen die einzelnen Retry-Strategien ein und führen einen Retry an den gegebenen Parametern aus.
Sobald eine Strategie keinen Retry mehr durchführen kann, sei es wegen Überschreitung des maximalen Retry-Counts oder
weil die Transaktion erfolgreich war, werden die Ergebnisse im JSONL-Format an die Funktion in
\texttt{experimentFunctions} zurückgegeben, wo diese dann in die Logdatei geschrieben werden.

Separat ist das \texttt{analysis} Modul.
Diesem wird ein Dateipfad zu den Logs übergeben und wertet diese aus.
Die Ergebnisse gibt es daraufhin strukturiert aus.


\section{Fehlererkennung}
\label{sec:errorDetection}
% !TeX root = ../../main.tex

Da die Datenbank für jede Transaktion einen SQLSTATE-Code zurückgibt, müssen diese Rückgabewerte entsprechend
ausgewertet werden.
Die Verarbeitung der SQLSTATE-Codes erfolgt im Modul \texttt{retry} innerhalb der Funktion
\texttt{\_start\_transaktion}.
Dort wird die Transaktion in einem \texttt{try: ... except: ...}-Block gestartet.
Tritt während der Ausführung ein von \texttt{psycopg2} Datenbankfehler auf, wird dieser abgefangen und
entsprechend verarbeitet~\ref{lst:errorDetection}.

\begin{lstlisting}[style=pythoncode, caption={Fehlererkennung}, label={lst:errorDetection}]
def _start_transaction():
    try:
        ...
    except psycopg2.Error as e:
        res = "failed"
        pgcode = e.pgcode
def run_retry_with_jitter_delay():
    ...
    while data == "failed" and retries < retry_count:
        ...
\end{lstlisting}

Im Fehlerfall wird neben dem von der Datenbank zurückgegebenen \texttt{pgcode} auch die Variable \texttt{res} auf
den Wert \texttt{failed} gesetzt.
Dadurch wird die weitere Verarbeitung unabhängig vom konkreten SQLSTATE-Code gesteuert.
Der \texttt{pgcode} wird dabei für die spätere Auswertung erfasst, während die Variable \texttt{res} zur Bestimmung des
weiteren Ablaufs des Retry-Mechanismus verwendet wird.

Die eigentliche Entscheidung über einen erneuten Ausführungsversuch erfolgt anschließend in der Retry Funktion.
Dazu wird zunächst geprüft, ob die vorherige Ausführung fehlgeschlagen ist.
Gleichzeitig wird überprüft, ob die maximal zulässige Anzahl an Retries bereits erreicht wurde.
Solange beide Bedingungen erfüllt sind, wird die Transaktion erneut ausgeführt.
Nach jedem fehlgeschlagenen, außer dem initial Versuch wird der Retry-Zähler erhöht, sodass die Anzahl der
erneuten Ausführungen begrenzt bleibt.
Sobald die Transaktion erfolgreich ausgeführt wurde oder die maximale Anzahl an Retries erreicht ist, wird die
Schleife verlassen und das vorliegende Ergebnis zurückgegeben.
Die Begrenzung der Anzahl der Retries ist insbesondere deshalb relevant, da ein Fehler andernfalls zu einer
unbegrenzten Folge von Wiederholungsversuchen führen könnte.

Für die vorliegende Untersuchung ist außerdem zu beachten, dass auf eine dynamische Klassifikation der Fehler
anhand der SQLSTATE-Codes verzichtet wurde.
Die Fehler werden in den Experimenten gezielt erzeugt und dienen jeweils der Untersuchung einer bestimmten
Fehlerklasse.
Stattdessen wird der zurückgegebene SQLSTATE-Code lediglich erfasst.
Im Rahmen der Versuchsauswertung wurde zusätzlich überprüft, dass bei der Durchführung der Experimente
ausschließlich die jeweils vorgesehenen Fehlerklassen auftraten.
Dadurch kann sichergestellt werden, dass die gemessenen Ergebnisse nicht durch unerwartete Fehlerarten beeinflusst
werden.

\section{Implementierung der Experimente}
\label{sec:experimentImpl}
% !TeX root = ../../main.tex

Die spezifischen Datenbankbefehle werden in dem Modul \texttt{database} definiert.
Dort steht die Class DB, über dessen Konstruktor die Datenbankverbindung aufgebaut wird.
Zusätzlich befinden sich Befehle für die Verwaltung in dieser Klasse, beispielsweise \texttt{initialize()} was erst
das Schema, die Tabellen für die einzelnen Fehlerklassen erstellt und daraufhin die Tabellen befüllt.
Befehle für die Einstellungen für die Serialisation und den Lock-Timeout befinden sich ebenfalls in dieser Klasse.

In den einzelnen \texttt{service\_....py}-Modulen werden die Transaktionen für die jeweiligen Fehlerklassen
implementiert.
Der Codeausschnitt~\ref{lst:deadlockImpl} zeigt beispielhaft die Implementierung eines Deadlocks.
Hierfür werden zwei Transaktionen verwendet, die jeweils ein Zeile der Tabelle \texttt{deadlocks} sperren und
anschließend versuchen, auf die von der anderen Transaktion gesperrte Zeile zuzugreifen.
Über die \texttt{client\_id} wird mit dem Modulo-Operator entschieden, welche der beiden Transaktionen ein Client
ausführen soll.

\begin{lstlisting}[language=Python, caption={Implementierung eines Deadlocks}, label={lst:deadlockImpl}]
def _tx_a(db):
    cur = db.conn.cursor()
    cur.execute("SELECT salary FROM deadlocks WHERE id = 1 FOR UPDATE;")
    sleep(0.05)
    cur.execute("UPDATE deadlocks SET salary = 100 WHERE id = 2;")
def _tx_b(db):
    cur = db.conn.cursor()
    cur.execute("SELECT salary FROM deadlocks WHERE id = 2 FOR UPDATE;")
    sleep(0.05)
    cur.execute("UPDATE deadlocks SET salary = 200 WHERE id = 1;")
def run_transaction_deadlocks(client_id, db):
    if client_id % 2 == 0:
        _tx_a(db)
    else:
        _tx_b(db)
\end{lstlisting}

Bei einer geraden \texttt{client\_id} wird somit die Funktion \texttt{\_tx\_a} ausgeführt, während bei einer
ungeraden \texttt{client\_id} die Transaktion in \texttt{\_tx\_b} zum Einsatz kommt.
Dadurch werden die Clients aufgeteilt und erzeugen einen Deadlock.
Die Funktion \texttt{\_tx\_a} sperrt zunächst den Datensatz mit der \texttt{id = 1}  und versucht anschließend den
Datensatz mit der \texttt{id = 2} zu aktualisieren.
\texttt{\_tx\_b} führt die gleichen Operationen in umgekehrter Reihenfolge aus.

Das Programm wird mit der \texttt{main}-Funktion gestartet.
Diese ruft nacheinander die Funktionen aus dem Modul \texttt{experimentRunner} auf.
Wie in Abschnitt~\ref{sec:softwareDesign} erwähnt, dienen die Funktionen in \texttt{experimentRunner}
ausschließlich als namensgebende Funktionen.
Dadurch wird abgeleitet und klar zugeordnet, welche Fehlerklassen erzeugt und durch die Experimente getestet werden
sollen.
Die Runner-Funktionen starten die Funktion \texttt{start()} in \texttt{experimentFunctions}.
Diese aufgerufene Funktion liest, über den in Coodeausschnitt~\ref{lst:caller-name} gezeigten Code, den Namen von der
aufrufenden Funktion aus.
Der Name wird bis zum Ende mit übergeben, damit die Funktionen wissen, welche Fehlerklassen behandelt und wohin die
Logeinträge geschrieben werden sollen.

\begin{lstlisting}[language=Python, caption={Ermittlung des Namens der aufrufenden Funktion}, label={lst:caller-name}]
	caller_frame = inspect.currentframe().f_back
	caller_name = caller_frame.f_code.co_name
\end{lstlisting}

\subsection{Modul experimentFunctions}
\label{subsec:experimentFunctionsImpl}
% !TeX root = ../../../main.tex

Des Weiteren werden dort Funktionen ausgeführt, um die Datenbank vorzubereiten.
Denn für jede Fehlerklasse wird die Datenbank neu aufgesetzt, damit alle zu gleichen Bedingungen starten.
Dafür wird die Datenbank vor Beginn der Experimente einem \emph{WARM-UP} unterzogen.
Dabei werden relevante Daten in die Caches der Datenbank geladen, wodurch nachfolgende Zugriffe schneller
verarbeitet werden können.
Damit die Ergebnisse durch vorherige Experimente nicht verfälscht werden, wird nach jedem Durchlauf eines
Experiments ein \emph{soft reset} der Datenbank durchgeführt.
Während dieser Zeit soll sich die Datenbank von Konflikten befreien, Locks lösen und wieder in den Idle Zustand gehen.
Dafür werden zunächst alle noch vorbereiteten Transaktionen der Datenbank mittels \texttt{ROLLBACK PREPARED}
zurückgesetzt.
Anschließend wird durch \texttt{sleep(2)} eine Wartezeit von zwei Sekunden eingehalten, damit sich die Datenbank
vollständig stabilisieren kann.
Nachdem die Datenbank vorbereitet ist, werden die Experimente ausgeführt.
Dafür wurden vier Experimente implementiert.
\texttt{experiment\_0} dient dabei als Baseline-Experiment und bildet den Ausgangspunkt für die weiteren Experimente.
Es beschreibt den Zustand des Systems ohne den Einsatz eines Retry-Mechanismus und ermöglicht somit einen Vergleich
mit den nachfolgenden Experimenten.
Dies wird für die fünf Concurrency-Werte getestet die in Ausschnitt~\ref{lst:params} gezeigt sind.
Das nächste Experiment, \texttt{experiment\_1}, soll den Einfluss von Retry ohne Delay testen.
Dadurch soll gezeigt werden, ob Retry ohne Delay sinnvoll ist und positive Ergebnisse liefert.
Hier ist insbesondere der \texttt{RETRY\_COUNT} Parameter, im folgenden nur noch Count genannt, wichtig, da dieser
anzeigt, wie sich das Verhalten ändert, wenn mehr Retries erlaubt sind.
Der Count von null wird in diesem Fall nicht beachtet, da dieser redundant ist.
In \texttt{experiment\_2} soll der Einfluss von Delay getestet werden.
Dafür wird der Count auf einen festen Wert gesetzt, hier Index 2 also Count = 3, damit dieser die Ergebnisse nicht
verfälscht und die \emph{immediate Retry} Strategie gewählt, weil die den Einfluss von Delay am besten zeigt, denn
der ist weder zufällig noch anderweitig veränderbar während der Strategie.
Der Delay ist nur in seinem Initialwert variabel, wobei auch der Delay von null nicht beachtet wird.
Als Letztes führt \texttt{experiment\_3} einen Vergleich der Strategien durch.
Es werden die beiden Strategien \emph{Retry with Delay} und \emph{Retry with backoff + jitter Delay}, mit einem festen
initial Delay und variabler Anzahl an Retries getestet.
Der initiale Delay ist auf 0,01 Sekunden festgelegt.
Dadurch lässt sich ein Vergleich zwischen allen Strategien ziehen, da bei allen die Anzahl der Retries variiert
wird.
Für die beiden Strategien, die einen Delay verwenden, wird mit 0,01 Sekunden ein bewusst geringer Delay gewählt,
sodass dieser keinen wesentlichen Einfluss auf die Vergleichbarkeit der Ergebnisse hat.
Da alle Experimente mit allen Concurrency-Werten durchgeführt werden, lässt sich in der Analyse herausfinden, ab
welcher Anzahl an gleichzeitigen Anfragen eine Retry Strategie
kippt.

Die Experimente rufen die Funktion \texttt{\_start\_experiment()} auf und übergeben die Indizes der Parameter.
Die aufgerufene Funktion liest die Werte der Parameter über die Funktion \texttt{\_get\_params()} anhand
ihrer Indizes aus.
Zusätzlich wird anhand der Parameter der \texttt{run\_key} erzeugt, der für die Sortierung in der Analyse benötigt
wird.
Die Parameter und der Schlüssel werden an \texttt{\_start\_experiment()} zurückgegeben.
Nachdem die Parameter und der Schlüssel gesetzt wurden, wird der erste Logeintrag für das Experiment erzeugt.
Die \emph{run\_meta} beschreibt genauer, worum es sich in dem Experiment handelt und wann es gestartet wurde.
Daraufhin wird \texttt{\_run\_batch()} aufgerufen, welche den Log-path anhand des caller\_name setzt und mit Threads
die Anfragen an die Datenbank startet.
Für die parallele Ausführung der Clients wird ein \texttt{ThreadPoolExecutor} mit der durch den Concurrency-Wert
angegebenen Anzahl an Threads verwendet.
Für jede Aufgabe wird ein \texttt{Future}-Objekt erzeugt und in einer Liste gespeichert.
Nachdem die Aufgaben abgeschlossen sind, werden die zugehörigen Ergebnisse über \texttt{as\_completed} nacheinander
verarbeitet.
Dadurch erfolgt das Schreiben in die Logdatei über den aufrufenden Thread, wodurch gleichzeitige Zugriffe auf die
Logdatei durch die Worker-Threads vermieden werden.
Für jeden ausgeführten Vorgang werden die Parameter \texttt{type}, \texttt{run\_id}, \texttt{iteration},
\texttt{retries}, \texttt{attempts\_ms}, \texttt{total\_ms}, \texttt{status}, \texttt{pg\_code}, \texttt{tx\_start}
und \texttt{tx\_finish} erfasst und im JSONL-Format in der Logdatei gespeichert.
Die Iteration wird nur wegen der Information mit geloggt, hat aber keinen weiteren Nutzen.

Im Client wird anhand der übergebenen Strategie entschieden, welche Strategie aufgerufen wird.
Das geschieht durch \texttt{match ... : case ... :} und in den einzelnen Cases wird der Rückgabewert der Strategie
gespeichert.
Die Datenbankverbindung wird ebenfalls in dieser Funktion aufgebaut, da jeder Client eine eigene Verbindung benötigt.
Die Ergebnisse werden dann an die zuvor erwähnte Funktion zurückgegeben.

\subsection{Modul retry}
\label{subsec:retryImpl}
% !TeX root = ../../../main.tex

Im Modul \texttt{retry} befinden sich die Implementierungen der einzelnen Retry Strategien.
Allen Strategien liegt dabei der ein gemeinsamer Ablauf zugrunde.
Zu Beginn werden in jeder Strategie die Variable \texttt{retries} und das Array \texttt{attempt\_times}
initialisiert, die sowohl die Anzahl an gebrauchten Retries, als auch die Ausführungsdauer der einzelnen
Transaktionsversuche speichern.
Danach wird ein Timestamp gesetzt, der den Beginn des gesamten Versuchs markiert.

Als Nächstes wird die in Codeausschnitt~\ref{lst:errorDetection} beschriebene Funktion aufgerufen.
Diese erkennt anhand des Parameters \texttt{caller\_name} welche Fehlerklasse erzeugt werden soll und ruft den
Service, der diese implementiert in einem \texttt{try:}-Block auf.
Als Rückgabewerte werden der Status der Transaktion, ihre Ausführungsdauer und der von PostgreSQL zurückgegebene
SQLSTATE-Code erfasst.
Die Dauer wird zusätzlich an das \texttt{attempt\_time}-Array angehängt.

Daraufhin werden die spezifischen Retry Mechanismen implementiert.
Beim Basisexperiment \texttt{run\_without\_retry} entfällt dieser Schritt, da die Transaktion lediglich einmal
ausgeführt wird.
Für dieses Experiment wird die Variable \texttt{retries} auf \texttt{None} gesetzt.
Dadurch wird verhindert, dass das Basisexperiment bei der späteren Auswertung der \emph{Retry Distribution} als
Experiment mit null Retries interpretiert wird.
Der Wert \texttt{None} wird bei der Analyse als \texttt{NaN} dargestellt und entsprechend separat ausgewertet.

Bei den Retry Mechanismen wird der Retry so lange durchgeführt, wie die maximale Anzahl an Retries es
zulässt, oder bis eine Transaktion erfolgreich war.
Dies geschieht mit einer \texttt{while}-Schleife in der als Erstes der Befehl \texttt{db.rollback()} durchgeführt,
damit sich die Verbindung wieder in einem fehlerfreien Zustand befindet.
Dann wird die Transaktion wiederholt und die zurückgegeben Werte werden gespeichert.
Die Ausführungszeit des Versuchs wird dem Array angehängt und die Anzahl der gebrauchten Retries wird um eins erhöht.
Dieser Ablauf wird wiederholt, bis eine der beiden Bedingung der Schleife nicht mehr zu treffen.

Die Implementierungen der einzelnen Strategien unterscheiden sich lediglich darin, wie sie den Delay behandeln.
Während bei der Strategie \emph{immediate Retry}-Strategie der Delay nicht verwendet wird, da die Transaktion, wird bei
\emph{Retry after Delay} der Delay auf den Wert des Parameters gesetzt.
Die \emph{Retry with Backoff + Jitter}-Strategie wird hingegen berechnet, dem maximalen Delay für jeden erneuten
Versuch, indem der initiale Delay mit $2^{n}$ multipliziert wird, wobei n der Anzahl der bereits durchgeführten
Retries entspricht.
Anschließend wird zufällig ein Delay zwischen null und dem berechneten maximal Wert gewählt~\ref{lst:jitter}.
In Python wird hierfür die Funktion \texttt{random.uniform} verwendet, welche einen gleichverteilten Zufallswert
innerhalb des angegebenen Intervalls erzeugt.
Die exponentielle Erhöhung des maximalen Delays wird durch \texttt{2 ** retries} umgesetzt.
Mit \texttt{time.sleep(delay)} wartet das System den definierten Delay ab und führt die Transaktion danach erneut aus.

\begin{lstlisting}[style=pythoncode, caption={Berechnung des Delay mit Backoff + Jitter}, label={lst:jitter}]
delay = random.uniform(0, retry_delay * (2 ** retries))
\end{lstlisting}

Der weitere Ablauf folgt für alle Strategien dem zuvor beschriebenen Muster.
Nach Ablauf des jeweiligen Delays wird die Transaktion erneut gestartet und anschließend anhand ihres
Rückgabestatus geprüft, ob sie erfolgreich ausgeführt werden konnte oder ein weiterer Retry erforderlich ist.
Sobald die Schleife beendet wird, wird ein Timestamp für das Ende des gesamten Versuchs erfasst.
Aus der Differenz zwischen End- und Startzeitpunkt wird anschließend die Gesamtdauer \texttt{total\_ms} des Versuchs
in Millisekunden berechnet.
Hierfür wird die Zeitdifferenz mit 1000 multipliziert und anschließend mit \texttt{math.floor} auf eine gerade Zahl
abgerundet.
Abschließend werden die gesammelten Ergebnisse, einschließlich der Anzahl der durchgeführten Retries, der
Ausführungszeiten der einzelnen Versuche, der Gesamtdauer, des Status und des SQLSTATE-Codes, zurückgegeben.



\subsection{Modul analysis}
\label{subsec:analysisImpl}
% !TeX root = ../../../main.tex

Das \texttt{analysis}-Modul enthält die Funktionen, die die Logdateien mit pandas auswerten.
Die Funktion \texttt{start\_analysis()} wird durch die \texttt{main} aufgerufen startet die einzelnen Analyseschritte.
Dafür wird die Logdatei eingelesen und in einem DataFrame gespeichert.

Durch die Funktion \texttt{\_print()} werden die Ergebnisse in die Konsole geschrieben.
Die Funktion in Codeausschnitt~\ref{lst:sortNumeric} implementiert eine numerische Sortierung.
Hierzu werden die ersten beiden Ebenen des MultiIndex betrachtet.
Sind die Indexwerte Strings, wird mithilfe eines regulären Ausdrucks der darin enthaltene numerische Wert
ausgelesen und in einen \texttt{float} umgewandelt.
Das \texttt{if hasattr} dient lediglich als Absicherung, dass nur in Strings die Zahl extrahiert wird, während bei
numerischen Werten direkt sortiert wird.
Dadurch werden die Ergebnisse in der numerisch korrekten Reihenfolge und nicht lexikografisch sortiert.
Wenn der Einfluss eines bestimmten Parameters betrachtet wird, steht dieser in dem ersten Level des Index, sodass die
Sortierung zunächst anhand dieses Parameters und anschließend anhand des zweiten Levels erfolgt.

\begin{lstlisting}[style=pythoncode, caption={Numerische Sortierung der Ergebnisse}, label={lst:sortNumeric}]
def _sort_numeric(data):
    return data.sort_index(
        level=[0, 1],
        key=lambda x: x.str.extract(r"([\d.]+)")[0].astype(float)
		if hasattr(x, 'str') else x
    )
\end{lstlisting}

Für die Auswertung stehen vier Funktionen zur Verfügung, die den Einfluss der Parameter Concurrency, Delay, Retry Count
und Strategie untersuchen.
Zusätzlich existiert eine Funktion, welche die Ergebnisse unabhängig von einem dieser Parameter auswertet.
Den vier parameterabhängigen Analysefunktionen wird jeweils der Parameter übergeben, dessen Einfluss auf die
Versuchsergebnisse untersucht werden soll.
Dieser Parameter bildet bei der anschließenden Auswertung mittels \texttt{groupby} die erste Ebene des Index und
besitzt somit die höchste Priorität bei der Sortierung und Auswertung der Ergebnisse.
Die Analysefunktionen rufen anschließend die einzelnen Funktionen zur Berechnung der in Abschnitt~\ref{sec:purpose}
beschriebenen Kenngrößen auf.
Dadurch kann dieselbe Berechnungslogik für unterschiedliche Parameter verwendet werden, ohne die eigentliche Berechnung
der Kenngrößen mehrfach implementieren zu müssen.

Im Folgendem wird nur der Fall, ohne zusätzliche Parameter betrachtet, da sich lediglich das Auslesen und die
Gruppierung nach diesem ändert.
Des Weiteren enthält jede Funktion, außer \texttt{\_throughput} den Parameter \texttt{pstatus}, welcher entweder auf
\texttt{success}, \texttt{failed} oder mit \texttt{""} leer gesetzt werden kann.
Dadurch kann das DataFrame wahlweise mit allen Einträgen, ausschließlich mit erfolgreichen Einträgen oder
ausschließlich mit fehlgeschlagenen Einträgen ausgegeben werden.

Die \texttt{\_rates()}-Funktion berechnet die Erfolgs- und Fehlerrate der Experimente in Prozent.
Hierfür wird das DataFrame zunächst nach dem Status gruppiert und die Anzahl der Einträge je Status gezählt und
gespeichert.
Die jeweilige Rate ergibt sich anschließend aus der Division von der Anzahl der Einträge und der Gesamtanzahl der
Einträge.
Das wird am Schluss mit 100 multipliziert, damit ein Prozentwert ausgegeben wird und nicht die Dezimalzahl.

Für die \texttt{\_retry\_distribution} wird das DataFrame nach \texttt{retries} und dem Status gruppiert.
Durch \texttt{dropna=False} werden \texttt{None}- beziehungsweise \texttt{NaN}-Werte berücksichtigt.
Dabei zählt \texttt{size()} die Einträge und \texttt{unstack()} wandelt die zweite Indexebene in Spalten um, während
\texttt{fill\_value=0} dafür sorgt, dass nicht vorhandene Kombinationen mit 0 dargestellt werden.

Bei der Funktion \texttt{\_execution\_comparison} werden Ausführungszeiten verglichen.
Dafür wird das DataFrame nach dem Status gruppiert und für jede Gruppe der Durchschnitt mit \texttt{mean()} von der
\texttt{total\_ms} berechnet.
Die Quartile in \texttt{\_quartiles} berechnen sich analog dazu, außer dass \texttt{quantile([0.50, 0.95, 0.99, 0.999])}
anstelle von \texttt{mean} verwendet wurde.

Zur Bestimmung des durch den Retry-Mechanismus entstehenden Overheads werden zunächst die einzelnen Ausführungszeiten
der Transaktionsversuche betrachtet.
Es wird aus \texttt{attempts\_ms} jeweils der erste Eintrag entfernt, da dieser die ursprüngliche Ausführung der
Transaktion und somit keinen Retry darstellt.
Aus den verbleibenden Ausführungszeiten wird anschließend mit \texttt{min()} die kürzeste Ausführungszeit bestimmt und
in \texttt{best\_ms} gespeichert.
Der Overhead ergibt sich schließlich aus der Differenz zwischen der gesamten Ausführungszeit \texttt{total\_ms} und der
kürzesten gemessenen Ausführungszeit.
Anschließend wird aus dem \texttt{overhead} der Durchschnitt berechnet.
Dadurch wird der zusätzliche Zeitaufwand bestimmt, der über die reine Ausführungszeit des schnellsten
Transaktionsversuchs hinausgeht.

Für die Berechnung des Durchsatzes wird zunächst aus der \texttt{run\_id} die jeweilige Experiment-ID extrahiert und in
der Spalte \texttt{e} gespeichert.
Anschließend werden die Einträge anhand dieser Experiment-ID gruppiert.
Für jedes Experiment werden dabei der früheste Startzeitpunkt und der späteste Endzeitpunkt aller Transaktionen
bestimmt.
Die Differenz zwischen diesen beiden Zeitpunkten ergibt die gesamte Laufzeit des jeweiligen Experiments.
Zusätzlich werden ausschließlich die erfolgreich abgeschlossenen Transaktionen betrachtet und pro Experiment gezählt.
Die Anzahl der erfolgreichen Transaktionen wird anschließend mit der zuvor bestimmten Laufzeit ins Verhältnis gesetzt.
Der resultierende Wert \texttt{throughput} gibt somit an, wie viele erfolgreiche Transaktionen pro Zeiteinheit
innerhalb des jeweiligen Experiments abgeschlossen wurden.

Durch diese Funktionen lassen sich die relevanten Kenngrößen berechnen und strukturiert nach den jeweiligen
Parametern darstellen.
Dadurch können die Ergebnisse gezielt miteinander verglichen und Rückschlüsse auf den Einfluss der
einzelnen Parameter gezogen werden.


\section{Konfiguration}
\label{sec:configuration}
% !TeX root = ../../main.tex

Der Codeausschnitt~\ref{lst:params} zeigt die Wertebereiche der für die Experimente verwendeten Parameter.
\texttt{WORKLOAD} und \texttt{ITERATION} bleiben über alle Experimente hinweg konstant und umfassen jeweils
500 auszuführende Transaktionen sowie fünf Wiederholungen.
Dadurch wird verhindert, dass Unterschiede in der Anzahl der ausgeführten Transaktionen die Messergebnisse
beeinflussen.
Unterschiede zwischen den Ergebnissen können somit auf die jeweils untersuchten Parameter zurückgeführt werden.

Die Parameter \texttt{CONCURRENCY}, \texttt{RETRY\_COUNT} und \texttt{RETRY\_DELAY} werden jeweils über mehrere
festgelegte Werte variiert.
Dadurch kann untersucht werden, wie sich beispielsweise eine steigende Anzahl paralleler Clients, eine höhere
Anzahl zulässiger Retries oder ein zunehmender Delay auf die jeweiligen Kenngrößen auswirkt.
Die Experimente werden dabei so konfiguriert, dass der jeweils untersuchte Parameter variiert wird, während die
übrigen Parameter konstant gehalten werden.
Dadurch lässt sich der Einfluss des einzelnen Parameters gezielter bestimmen.

Der Parameter \texttt{RETRY\_STRATEGY} ist aus Gründen der Optimierung ebenfalls in der \texttt{config.py} definiert.
Die Auswahl der tatsächlich verwendeten Strategie erfolgt jedoch innerhalb der Implementierung über eine
\texttt{match/case}-Verzweigung.
Dadurch können die verschiedenen Retry-Strategien mit derselben grundlegenden Logik ausgeführt und miteinander
verglichen werden.

Die gewählten Parameterbereiche stellen dabei einen Kompromiss zwischen dem Umfang der Experimente und der
Aussagekraft der Ergebnisse dar.
Durch die Variation mehrerer Werte können Trends innerhalb der untersuchten Kenngrößen erkannt werden, ohne dass
für jede mögliche Kombination aller Parameter ein eigenes Experiment durchgeführt werden muss.
Die fünf Wiederholungen pro Konfiguration dienen zusätzlich dazu, Schwankungen einzelner Messungen zu reduzieren
und eine stabilere Grundlage für die spätere Auswertung zu schaffen.

\begin{lstlisting}[style=pythoncode, caption={Übersicht der Parameter}, label={lst:params}]
RETRY_COUNT = [0, 1, 3, 5]
RETRY_DELAY = [0, 0.005, 0.01, 0.05, 0.1]
RETRY_STRATEGY = [0, 1, 2, 3]

CONCURRENCY = [2, 5, 7, 10, 20]
WORKLOAD = [500]
ITERATIONS = 5
\end{lstlisting}

\section{Implementierungsentscheidungen}
\label{sec:decisions}
% !TeX root = ../../main.tex

Bei der Implementierung wurden verschiedene Entscheidungen getroffen, um den Versuchsaufbau möglichst einfach und
reproduzierbar zu halten.
Gleichzeitig sollte die Implementierung möglichst nah an den tatsächlich untersuchten Datenbank- und
Retry-Mechanismen bleiben.
Die wichtigsten Entscheidungen werden im Folgenden kurz begründet.

Für die Umsetzung wurde Python verwendet, da sich damit sowohl die Kommunikation mit PostgreSQL als auch die
parallele Ausführung der Transaktionen und die spätere Analyse der Versuchsdaten umsetzen lassen.
Für die Datenbankanbindung wird \texttt{psycopg2} eingesetzt, damit ein direkter Zugriff auf PostgreSQL
möglich ist.
Dadurch können SQL-Anweisungen und Transaktionen gezielt gesteuert und auftretende Datenbankfehler direkt
verarbeitet werden.
Zudem handelt es sich bei \texttt{psycopg2} um eine etablierte und weit verbreitete Bibliothek für die Anbindung
von PostgreSQL aus Python.
Als Datenbanksystem wird PostgreSQL verwendet, da die für die Experimente benötigten Transaktionsmechanismen und
Sperrverfahren unterstützt werden.
Insbesondere die Erkennung von Fehlern und die Rückgabe entsprechender SQLSTATE-Codes sind für die Untersuchung der
Fehlerklassen relevant.
Die Datenbank wird lokal betrieben, um zusätzliche Einflüsse durch Netzwerkverbindungen oder externe Systeme zu
vermeiden.
Dadurch kann die Testumgebung zudem einfacher konfiguriert und kontrolliert werden.

Für die Speicherung der Versuchsergebnisse wurde das JSONL-Format gewählt.
Jeder Transaktionsversuch kann dabei als eigener Datensatz gespeichert werden, wodurch auch bei einer großen Anzahl
an Messwerten eine einfache und strukturierte Speicherung möglich ist.
Die erzeugten Logdateien können anschließend direkt mit \texttt{pandas} eingelesen und weiterverarbeitet werden.
\texttt{pandas} wird für die Gruppierung der Versuchsdaten sowie die Berechnung der untersuchten Kenngrößen
eingesetzt.

Die Fehler werden gezielt erzeugt, um für jedes Experiment eine definierte Fehlerklasse untersuchen zu können.
Eine dynamische Klassifikation anhand der SQLSTATE-Codes ist daher für die Durchführung der Experimente nicht
notwendig.
Die Codes werden dennoch protokolliert, um die tatsächlich aufgetretenen Fehler überprüfen und die Ergebnisse
entsprechend nachvollziehen zu können.

Für die Untersuchung des Retry-Verhaltens wurden die Strategien \emph{Immediate Retry}, \emph{Retry after Delay} und
\emph{Retry with Backoff + Jitter} gewählt.
Diese decken unterschiedliche Ansätze der Wiederholungsversuche ab.
Während beim \emph{Immediate Retry} unmittelbar ein neuer Versuch gestartet wird, verwendet \emph{Retry after Delay}
eine konstante Wartezeit.
\emph{Retry with Backoff + Jitter} erweitert diesen Ansatz um einen exponentiell steigenden maximalen Delay und
eine zufällige Komponente.
Dadurch können die Auswirkungen verschiedener Verzögerungsstrategien auf die untersuchten Kenngrößen miteinander
verglichen werden.

Die parallele Ausführung der Clients erfolgt über den \texttt{ThreadPoolExecutor}.
Dadurch können mehrere Transaktionen gleichzeitig ausgeführt und die für die Experimente benötigte Concurrency
realisiert werden.
Der nächste Versuch wird erst gestartet, nachdem der vorherige fehlgeschlagen ist und gegebenenfalls der
entsprechende Delay abgewartet wurde.
